% ============================================================
%  Technická zpráva – šablona
%  Robotika / Robolab, FEE ČVUT (nebo jiný ústav)
%  Formát: A4, jednostranný tisk
% ============================================================
\documentclass[12pt, a4paper, oneside]{article}

% ── Kódování a jazyk ─────────────────────────────────────────
\usepackage[utf8]{inputenc}
\usepackage[T1]{fontenc}
\usepackage[czech]{babel}

% ── Písmo ────────────────────────────────────────────────────
\usepackage{lmodern}          % Moderní Latin Modern

% ── Geometrie stránky ────────────────────────────────────────
\usepackage[
  a4paper,
  top=25mm, bottom=25mm,
  left=30mm, right=25mm
]{geometry}

% ── Matematika ───────────────────────────────────────────────
\usepackage{amsmath}
\usepackage{amssymb}
\usepackage{amsthm}

% ── Grafika a tabulky ────────────────────────────────────────
\usepackage{graphicx}
\usepackage{float}            % Přesné umístění obrázků [H]
\usepackage{booktabs}         % Profesionální tabulky
\usepackage{caption}
\usepackage{subcaption}       % Vedlejší obrázky / tabulky

% ── Kód a pseudokód ─────────────────────────────────────────
\usepackage{listings}
\usepackage{xcolor}

\lstset{
  basicstyle=\ttfamily\small,
  keywordstyle=\color{blue}\bfseries,
  commentstyle=\color{gray}\itshape,
  stringstyle=\color{teal},
  numbers=left,
  numberstyle=\tiny\color{gray},
  stepnumber=1,
  numbersep=6pt,
  breaklines=true,
  frame=single,
  captionpos=b,
  showstringspaces=false,
}

% ── Algoritmy ────────────────────────────────────────────────
\usepackage[ruled, vlined, linesnumbered]{algorithm2e}
\SetKwInput{KwIn}{Vstup}
\SetKwInput{KwOut}{Výstup}

% ── Hypertextové odkazy ──────────────────────────────────────
\usepackage[
  colorlinks=true,
  linkcolor=black,
  citecolor=black,
  urlcolor=blue
]{hyperref}

% ── Záhlaví a zápatí ────────────────────────────────────────
\usepackage{fancyhdr}
\pagestyle{fancy}
\fancyhf{}
\fancyhead[L]{\small\leftmark}
\fancyhead[R]{\small Technická zpráva}
\fancyfoot[C]{\thepage}
\renewcommand{\headrulewidth}{0.4pt}

% ── Řádkování ────────────────────────────────────────────────
\usepackage{setspace}
\onehalfspacing

% ── Bibliografie ─────────────────────────────────────────────
\usepackage[
  backend=biber,
  style=numeric,
  sorting=none
]{biblatex}
\addbibresource{references.bib}   % Soubor s literaturou

% ── Pomocné balíčky ──────────────────────────────────────────
\usepackage{microtype}        % Lepší zalamování řádků
\usepackage{parskip}          % Mezera mezi odstavci místo odsazení
\usepackage{enumitem}         % Přizpůsobení seznamů

% ============================================================
%  VLASTNÍ PŘÍKAZY – vyplňte před tiskem
% ============================================================
\newcommand{\NazevPrace}{Autonomní robot Ferenc}
\newcommand{\Autorka}{Jakub Kecek}
\newcommand{\AutorB}{Václav Kučera}
\newcommand{\AutorC}{Ondřej Vosáhlo}
\newcommand{\DatumOdevzdani}{\today}
\newcommand{\Ustav}{Fakulta elektrotechnická, ČVUT v Praze}

% ============================================================
\begin{document}

% ── Titulní strana ───────────────────────────────────────────
\begin{titlepage}
  \centering
  \vspace*{2cm}

  {\Large \Ustav \par}
  \vspace{0.5cm}
  {\large \Predmet \par}

  \vspace{3cm}

  {\Huge \bfseries \NazevPrace \par}

  \vspace{2cm}

  {\large
    \begin{tabular}{ll}
      \textbf{Autoři:}   & \Autorka \\
                         & \AutorB  \\
                         & \AutorC  \\[6pt]
      \textbf{Datum:}    & \DatumOdevzdani \\
    \end{tabular}
  }

  \vfill

  {\small Zpráva vysázena pomocí \LaTeX{} na stránkách formátu A4.}
\end{titlepage}

% ── Obsah ────────────────────────────────────────────────────
\tableofcontents
\newpage

% ============================================================
%  1. POPIS PROBLÉMU
% ============================================================
\section{Popis řešeného problému}
\label{sec:popis}

% TODO: Stručně popište zadání úlohy. Co je cílem robota?
%       Co musí robot detekovat, jaké úkony provést a v jakém pořadí?

Robot \textit{Ferenc} má za úkol autonomně opustit garáž, lokalizovat míček
v prostoru, objíždět jej po předdefinované trajektorii a poté se vrátit
do výchozí pozice. Celý systém je implementován na platformě Turtlebot využitím RGB kamery a 3D hloubkového senzoru (point cloud).

% ============================================================
%  2. ROZBOR PROBLÉMU A NÁVRH ŘEŠENÍ (SPECIFIKACE)
% ============================================================
\section{Rozbor problému, návrh řešení a očekávaná funkčnost}
\label{sec:specifikace}

\subsection{Analýza vstupů a výstupů}

% TODO: Popište, jaké senzory jsou k dispozici a jaké výstupy je třeba
%       produkovat (příkazy rychlosti, navigační rozhodnutí apod.).

\subsection{Dekompozice úlohy}

Celou úlohu lze rozdělit do následujících kroků:

\begin{enumerate}[label=\arabic*.]
  \item Detekce výjezdu z garáže a vyjetí do volného prostoru.
  \item Lokalizace míčku pomocí barevné segmentace v obraze kamery.
  \item Přiblížení k míčku na požadovanou vzdálenost.
  \item Objezd míčku po hexagonální trajektorii.
  \item Návrat do výchozí pozice pomocí odometrické navigace.
\end{enumerate}

\subsection{Očekávaná funkčnost}

% TODO: Popište, jak by měl hotový systém fungovat za ideálních podmínek.

\subsection{Matematický popis a algoritmy}
\label{sec:matematika}

\subsubsection{Detekce míčku}

Detekce je prováděna v barevném prostoru HSV. Pixel $(u, v)$ je označen jako
součást míčku, platí-li:

\begin{equation}
  H_{\min} \le H(u,v) \le H_{\max}, \quad
  S(u,v) \ge S_{\min}, \quad
  V(u,v) \ge V_{\min}.
  \label{eq:hsv_threshold}
\end{equation}

Kruhovitost každého nalezeného konturu $C$ je definována jako:

\begin{equation}
  \kappa(C) = \frac{4\pi \cdot A(C)}{P(C)^2},
  \label{eq:circularity}
\end{equation}

kde $A(C)$ je plocha a $P(C)$ obvod konturu. Kontur je přijat jako míček,
splňuje-li $\kappa(C) > \kappa_{\min}$.

\subsubsection{Odhad vzdálenosti}

Vzdálenost k míčku je odhadována ze středních hodnot hloubkového senzoru
v okolí středu detekovaného konturu:

\begin{equation}
  d = \frac{1}{|W|} \sum_{(u,v)\in W} z(u,v),
  \label{eq:depth_estimate}
\end{equation}

kde $W$ je okno velikosti $(2r+1) \times (2r+1)$ se středem v $(u_0, v_0)$
a $r \in \{1, 2\}$ závisí na poloměru detekovaného míčku v pixelech.

\subsubsection{PID regulátor otáčení}

Úhlová odchylka $e_\theta(t)$ je regulována PID regulátorem:

\begin{equation}
  u(t) = K_p\, e_\theta(t)
         + K_i \int_0^t e_\theta(\tau)\,\mathrm{d}\tau
         + K_d \frac{\mathrm{d}e_\theta}{\mathrm{d}t}.
  \label{eq:pid}
\end{equation}

Použité konstanty jsou $K_p = 2{,}1$, $K_i = 0{,}067$, $K_d = 0{,}18$.

\subsubsection{Hexagonální trajektorie}

Střed míčku v souřadnicích odometrie je označen jako $\mathbf{b} = (b_x, b_y)$.
$i$-tý uzel hexagonu (pro $i = 0, 1, \ldots, 5$) leží na kružnici
o poloměru $r_h = d_{\text{stop}} + r_{\text{míček}}$:

\begin{equation}
  \mathbf{w}_i = \mathbf{b} +
    r_h \begin{pmatrix} \cos\!\left(\dfrac{i\pi}{3}\right) \\[4pt]
                        \sin\!\left(\dfrac{i\pi}{3}\right) \end{pmatrix}.
  \label{eq:hexagon}
\end{equation}

% ============================================================
%  3. ZMĚNY SPECIFIKACE
% ============================================================
\section{Změny specifikace v průběhu řešení}
\label{sec:zmeny}

% TODO: Uveďte, co se od původního plánu změnilo, proč a jaký to mělo dopad.

\begin{itemize}
  \item \textbf{Změna~1:} \ldots
  \item \textbf{Změna~2:} \ldots
\end{itemize}

% ============================================================
%  4. IMPLEMENTACE
% ============================================================
\section{Implementace a způsob spuštění}
\label{sec:implementace}

\subsection{Struktura projektu}

Projekt je rozdělen do tří modulů:

\begin{description}
  \item[\texttt{ferenc.py}] Hlavní řídicí logika robota – třída \texttt{Ferenc}.
  \item[\texttt{visuals.py}] Detekce míčku a obdélníkových značek.
  \item[\texttt{image\_processing.py}] Zpracování hloubkového obrazu a detekce volného prostoru.
\end{description}

\subsection{Způsob spuštění}

% TODO: Doplňte konkrétní příkazy nebo kroky pro spuštění.

Robot se spustí příkazem:

\begin{lstlisting}[language=bash, caption={Spuštění robota.}]
python3 ferenc.py
\end{lstlisting}

Vyžaduje prostředí \texttt{robolab\_turtlebot}, OpenCV a NumPy.

% ============================================================
%  5. EXPERIMENTÁLNÍ VÝSLEDKY
% ============================================================
\section{Experimentální výsledky}
\label{sec:vysledky}

% TODO: Doplňte tabulku s výsledky pro jednotlivé pokusy / scénáře.
%       Uveďte podmínky experimentu (osvětlení, poloha míčku apod.).

\begin{table}[H]
  \centering
  \caption{Výsledky experimentů – přiblížení k míčku a objezd.}
  \label{tab:vysledky}
  \begin{tabular}{ccccc}
    \toprule
    \textbf{Pokus} & \textbf{Vzdál. míčku [m]} & \textbf{Chyba zastavení [m]}
                   & \textbf{Objezd dokončen} & \textbf{Návrat do garáže} \\
    \midrule
    1 & 1{,}20 & 0{,}03 & ano & ano \\
    2 & 0{,}85 & 0{,}05 & ano & ne  \\
    3 & 1{,}50 & 0{,}02 & ano & ano \\
    % Doplňte další řádky dle potřeby
    \bottomrule
  \end{tabular}
\end{table}

% ============================================================
%  6. DISKUSE VÝSLEDKŮ
% ============================================================
\section{Diskuse výsledků}
\label{sec:diskuse}

% TODO: Zhodnoťte, co fungovalo dobře, kde nastávaly problémy.
%       Co byste zlepšili s více čase a prostředky?

\subsection{Silné stránky řešení}

\subsection{Omezení a nedostatky}

\subsection{Návrhy na zlepšení}

Při větším objemu zdrojů by bylo vhodné:

\begin{itemize}
  \item Nahradit odometrický návrat lokalizací na základě simultánní lokalizace
        a mapování (SLAM).
  \item Implementovat robustnější detekci míčku pomocí hluboké neuronové sítě.
  \item Přidat zpracování výjimek při kolizi s neočekávanou překážkou.
\end{itemize}

% ============================================================
%  7. NÁVRH ÚLOHY PRO PŘÍŠTÍ ROK
% ============================================================
\section{Návrh úlohy pro příští rok}
\label{sec:navrh}

% TODO: Navrhněte rozšíření nebo alternativní zadání pro budoucí cvičení.

% ============================================================
%  8. LITERATURA
% ============================================================
\printbibliography[heading=bibintoc, title={Literatura}]

% ============================================================
%  9. PŘÍLOHY
% ============================================================
\appendix

\section{Příloha A – Úplný zdrojový kód}
\label{app:kod}

% Příklad vložení části kódu přímo do zprávy.
\begin{lstlisting}[language=Python, caption={Výpočet průměrné hloubky (zkráceno).}]
def get_depth(turtle, center_x, center_y, radius):
    if radius < 2:
        return None
    pc = turtle.get_point_cloud()
    r = 1 if radius < 16 else 2
    depths = [
        float(pc[center_y + i][center_x + j][2])
        for i in range(-r, r + 1)
        for j in range(-r, r + 1)
        if pc[center_y + i][center_x + j][2] is not None
        and float(pc[center_y + i][center_x + j][2]) > 0.1
    ]
    return sum(depths) / len(depths) if depths else None
\end{lstlisting}

\section{Příloha B – Fotografie z experimentu}
\label{app:foto}

% \begin{figure}[H]
%   \centering
%   \includegraphics[width=0.7\textwidth]{foto_experiment.jpg}
%   \caption{Robot Ferenc při objezdu míčku.}
%   \label{fig:foto}
% \end{figure}

\end{document}
